\documentclass[11pt]{article}

\usepackage[a4paper,margin=22mm]{geometry}
\usepackage{fontspec}
\usepackage{xeCJK}
\usepackage{amsmath,amssymb,mathtools,bm}
\usepackage{booktabs}
\usepackage{enumitem}
\usepackage{xcolor}
\usepackage[hidelinks]{hyperref}
\usepackage{fancyhdr}

\setmainfont{Helvetica Neue}
\setsansfont{Helvetica Neue}
\setmonofont{Menlo}
\setCJKmainfont{Songti SC}
\setCJKsansfont{PingFang SC}
\setCJKmonofont{PingFang SC}

\hypersetup{
  pdftitle={Cyrus CS231n Guided Notes},
  pdfauthor={Cyrus Learning Manager},
  pdfsubject={Stanford CS231n guided study notes},
  pdfkeywords={computer vision, deep learning, CNN, spatial intelligence}
}

\pagestyle{fancy}
\fancyhf{}
\lhead{\small Cyrus CS231n Guided Notes}
\rhead{\small \thepage}
\renewcommand{\headrulewidth}{0.4pt}

\setlength{\parindent}{0pt}
\setlength{\parskip}{0.55em}
\setlist[itemize]{leftmargin=1.4em,itemsep=0.2em,topsep=0.2em}
\setlist[enumerate]{leftmargin=1.6em,itemsep=0.2em,topsep=0.2em}
\sloppy

\newcommand{\R}{\mathbb{R}}
\newcommand{\softmax}{\operatorname{softmax}}
\newcommand{\ReLU}{\operatorname{ReLU}}
\newcommand{\BN}{\operatorname{BN}}
\newcommand{\E}{\mathbb{E}}
\newcommand{\KL}{D_{\mathrm{KL}}}
\newcommand{\goodnotes}[1]{\textbf{GoodNotes 任务：}#1}

\begin{document}

\begin{titlepage}
  \centering
  {\sffamily\bfseries\Huge Cyrus CS231n 导学讲义\par}
  \vspace{0.8em}
  {\Large Deep Learning for Computer Vision\par}
  \vspace{1.2em}
  {\large 按“李飞飞 CS231n / CS23DN”学习需求整理的中文路线版\par}
  \vfill
  \begin{tabular}{ll}
    官方课程 & \href{https://cs231n.stanford.edu/}{Stanford CS231n} \\
    官方 notes & \href{https://cs231n.github.io/}{cs231n.github.io} \\
    本地课程目录 & \texttt{docs/courses/cs231n/} \\
    学习方式 & 网页交互 + GoodNotes 推导 + Obsidian 概念图 + Notion 证据记录 \\
  \end{tabular}
  \vfill
  {\small 本文是原创导学笔记和公式索引，不替代官方课程页面、slides、assignments 或 notes。}
\end{titlepage}

\tableofcontents
\newpage

\section*{怎么用这份讲义}
\addcontentsline{toc}{section}{怎么用这份讲义}

CS231A 解决“相机和几何如何恢复 3D”；CS231n 解决“神经网络如何从图像中学习表示”。你的空间智能路线需要二者合起来：
\[
  \text{CS231A geometry} + \text{CS231n visual representation}
  \Rightarrow
  \text{modern spatial intelligence}.
\]

每一章按四件事学习：
\begin{enumerate}
  \item 概念：它解决什么视觉问题。
  \item 公式：损失函数、梯度、层结构或概率目标。
  \item GoodNotes：手抄核心公式并画一张结构图。
  \item 网页/Notion：记录掌握度、错题、作业输出和下一步。
\end{enumerate}

\section{课程地图}

\begin{center}
\begin{tabular}{p{0.08\linewidth}p{0.32\linewidth}p{0.50\linewidth}}
\toprule
编号 & 模块 & 学习目标 \\
\midrule
1 & Image Classification & 数据驱动分类、kNN、train/val/test。 \\
2 & Linear Classifiers & SVM、Softmax、正则化和概率解释。 \\
3 & Optimization and Backprop & SGD、链式法则、计算图反向传播。 \\
4 & Neural Networks & 多层网络、激活函数、loss、初始化。 \\
5 & CNNs & 卷积、池化、感受野、经典架构。 \\
6 & Training Tricks & BatchNorm、Dropout、优化器、迁移学习。 \\
7 & Dense Prediction & Detection、segmentation、IoU、NMS。 \\
8 & Sequence and Attention & RNN、LSTM、attention、Transformer。 \\
9 & Generative Models & Autoencoder、VAE、GAN、diffusion 入口。 \\
10 & Self-supervised Vision & Contrastive learning、CLIP/ViT 入口、空间智能连接。 \\
\bottomrule
\end{tabular}
\end{center}

\section{Image Classification}

\subsection{数据驱动视觉}

传统视觉会手工设计边缘、角点、纹理等特征；CS231n 的起点是数据驱动：给训练集
\[
  \{(x_i,y_i)\}_{i=1}^{N},
\]
学习一个函数
\[
  f_\theta(x)\rightarrow y.
\]
图像分类的目标是把输入图像 $x$ 映射到类别标签 $y$。

\subsection{kNN}

kNN 不学习显式参数，只存训练样本。给新图像 $x$，计算它和训练图像 $x_i$ 的距离。常见距离：
\[
  d_1(x,x_i)=\sum_p |x_p-x_{i,p}|,
\]
\[
  d_2(x,x_i)=\sqrt{\sum_p (x_p-x_{i,p})^2}.
\]
选择最近的 $k$ 个样本投票。它简单但测试慢，且像素距离并不等于语义相似。

\subsection{数据划分}

训练集用于拟合参数，验证集用于选超参数，测试集只做最后评估：
\[
  \text{train}\rightarrow \text{validation}\rightarrow \text{test}.
\]
空间智能里同理：训练数据、验证场景、最终验收路线不能混用。

\goodnotes{画出 train/val/test 三块数据，写下 L1、L2 距离，并解释为什么像素近不代表语义近。}

\section{Linear Classifiers}

\subsection{线性打分函数}

线性分类器把图像拉平成向量 $x\in\R^D$，输出每类分数：
\[
  s=f(x;W,b)=Wx+b,
\]
其中 $W\in\R^{C\times D}$，$b\in\R^C$。第 $j$ 类分数为 $s_j$。

\subsection{Multiclass SVM Loss}

若正确类别为 $y_i$，SVM 希望正确类分数比其他类至少高 $\Delta$：
\[
  L_i=
  \sum_{j\neq y_i}
  \max(0,s_j-s_{y_i}+\Delta).
\]
总损失加正则：
\[
  L=\frac{1}{N}\sum_i L_i+\lambda R(W),
\qquad
  R(W)=\sum_k W_k^2.
\]
直觉：分类器不仅要答对，还要答得有 margin。

\subsection{Softmax and Cross Entropy}

Softmax 把分数转成概率：
\[
  p(y=j\mid x)=\frac{e^{s_j}}{\sum_k e^{s_k}}.
\]
交叉熵损失：
\[
  L_i=-\log p(y_i\mid x_i)
  =
  -s_{y_i}+\log\sum_k e^{s_k}.
\]
SVM 强调 margin；Softmax 强调概率校准和最大似然。

\goodnotes{写出 $s=Wx+b$，画出一个三分类分数向量，分别用 SVM 和 Softmax 判断损失。}

\section{Optimization and Backpropagation}

\subsection{梯度下降}

训练就是最小化损失：
\[
  \theta^\star=\arg\min_\theta L(\theta).
\]
梯度下降更新：
\[
  \theta_{t+1}=\theta_t-\eta\nabla_\theta L(\theta_t),
\]
其中 $\eta$ 是 learning rate。SGD 用小批量估计梯度：
\[
  \nabla_\theta L(\theta)\approx
  \frac{1}{B}\sum_{i\in\mathcal B}\nabla_\theta L_i(\theta).
\]

\subsection{链式法则}

若
\[
  z=f(y),\qquad y=g(x),
\]
则
\[
  \frac{\partial z}{\partial x}
  =
  \frac{\partial z}{\partial y}
  \frac{\partial y}{\partial x}.
\]
反向传播就是在计算图上反复使用链式法则，把最终 loss 的梯度传回每个参数。

\subsection{反传的局部视角}

每个节点只需要知道两件事：
\begin{itemize}
  \item 前向：自己的输出；
  \item 反向：上游梯度乘本地导数。
\end{itemize}
例如乘法 $z=xy$：
\[
  \frac{\partial z}{\partial x}=y,\qquad
  \frac{\partial z}{\partial y}=x.
\]
ReLU：
\[
  \ReLU(x)=\max(0,x),
\qquad
  \frac{\partial \ReLU}{\partial x}=
  \begin{cases}
    1,&x>0,\\
    0,&x\le 0.
  \end{cases}
\]

\goodnotes{画一个 $x\rightarrow y\rightarrow L$ 的计算图，标出前向值和反向梯度。}

\section{Neural Networks}

\subsection{从线性到非线性}

线性分类器只能学习线性决策边界。多层网络堆叠 affine layer 和非线性激活：
\[
  h_1=\ReLU(W_1x+b_1),
\]
\[
  s=W_2h_1+b_2.
\]
多层网络可表示更复杂的视觉模式，但也更难训练。

\subsection{激活函数}

Sigmoid：
\[
  \sigma(x)=\frac{1}{1+e^{-x}}.
\]
Tanh：
\[
  \tanh(x)=\frac{e^x-e^{-x}}{e^x+e^{-x}}.
\]
ReLU：
\[
  \ReLU(x)=\max(0,x).
\]
现代视觉网络常用 ReLU/GELU/SiLU 类激活，因为梯度传播更稳定。

\subsection{初始化和正则化}

若初始化太小，信号逐层消失；太大，激活爆炸。常见初始化让方差在层间近似保持稳定，例如 He initialization：
\[
  W_{ij}\sim \mathcal N\left(0,\frac{2}{n_{\text{in}}}\right).
\]
正则化防止过拟合：
\[
  L_{\text{total}}=L_{\text{data}}+\lambda\|W\|_2^2.
\]

\section{Convolutional Neural Networks}

\subsection{卷积层}

卷积利用图像的局部性和平移共享。若输入空间尺寸为 $W$，卷积核尺寸为 $F$，padding 为 $P$，stride 为 $S$，输出尺寸：
\[
  W_{\text{out}}=\frac{W-F+2P}{S}+1.
\]
若卷积核数量为 $K$，输入通道为 $C$，参数量：
\[
  F\cdot F\cdot C\cdot K+K.
\]
共享参数让 CNN 比全连接层更适合图像。

\subsection{Pooling and Receptive Field}

Pooling 汇聚局部响应，降低空间尺寸并增加局部不变性。感受野表示输出单元能看到输入图像的范围。层数越深，感受野越大，语义越抽象。

\subsection{经典架构}

\begin{itemize}
  \item AlexNet：证明深度 CNN 在 ImageNet 上有效。
  \item VGG：使用堆叠 $3\times3$ 卷积，结构规整。
  \item GoogLeNet/Inception：多尺度分支与 $1\times1$ 卷积降维。
  \item ResNet：残差连接
  \[
    y=F(x)+x
  \]
  缓解深层网络优化困难。
\end{itemize}

\goodnotes{画一条 CNN 流水线：image $\rightarrow$ conv $\rightarrow$ ReLU $\rightarrow$ pool $\rightarrow$ classifier。}

\section{Training Tricks and Transfer Learning}

\subsection{Batch Normalization}

BatchNorm 对 mini-batch 内的激活标准化：
\[
  \mu_B=\frac{1}{m}\sum_i x_i,
\qquad
  \sigma_B^2=\frac{1}{m}\sum_i (x_i-\mu_B)^2,
\]
\[
  \hat x_i=\frac{x_i-\mu_B}{\sqrt{\sigma_B^2+\epsilon}},
\qquad
  y_i=\gamma \hat x_i+\beta.
\]
它改善优化稳定性，也带来一定正则效果。

\subsection{Dropout}

训练时随机丢弃神经元：
\[
  m_i\sim \operatorname{Bernoulli}(p),
\qquad
  \tilde h_i=\frac{m_i h_i}{p}.
\]
这迫使网络不依赖单个特征。

\subsection{Momentum and Adam}

Momentum：
\[
  v_{t+1}=\mu v_t-\eta\nabla_\theta L(\theta_t),
\qquad
  \theta_{t+1}=\theta_t+v_{t+1}.
\]
Adam 使用一阶矩和二阶矩估计：
\[
  m_t=\beta_1m_{t-1}+(1-\beta_1)g_t,
\]
\[
  v_t=\beta_2v_{t-1}+(1-\beta_2)g_t^2.
\]

\subsection{Transfer Learning}

若数据少，直接从头训练大模型不现实。迁移学习流程：
\[
  \text{pretrained backbone}\rightarrow \text{freeze or finetune}\rightarrow \text{task head}.
\]
对你的空间智能路线，ImageNet/CLIP/DINO/ViT backbone 可以作为检测、分割、深度、地图语义的表示层。

\section{Detection and Segmentation}

\subsection{Detection}

目标检测输出类别和边界框。框质量常用 IoU：
\[
  \operatorname{IoU}(A,B)=
  \frac{|A\cap B|}{|A\cup B|}.
\]
NMS 过滤重复框：
\[
  \text{keep high score box, remove boxes with IoU}>\tau.
\]
Two-stage 检测器先提 proposals，再分类和回归；one-stage 检测器直接密集预测。

\subsection{Segmentation}

语义分割给每个像素类别：
\[
  f_\theta(I)_{u,v}\rightarrow c.
\]
实例分割还区分同类不同物体。对于空间智能，segmentation 是把图像转成可用于地图、规划和语言 grounding 的结构化表示。

\subsection{Dense Prediction}

深度估计、光流、分割都属于 dense prediction：输出不是一个类别，而是每个像素一个值或向量。典型结构是 encoder-decoder：
\[
  I\rightarrow z\rightarrow \hat Y.
\]

\section{RNN, Attention, and Transformers}

\subsection{RNN}

序列模型维护隐藏状态：
\[
  h_t=\phi(W_{xh}x_t+W_{hh}h_{t-1}+b_h),
\]
\[
  y_t=W_{hy}h_t+b_y.
\]
RNN 可用于图像描述、视频理解和时序状态建模，但长序列训练困难。

\subsection{LSTM}

LSTM 用门控控制记忆：
\[
  f_t=\sigma(W_f[h_{t-1},x_t]+b_f),
\]
\[
  i_t=\sigma(W_i[h_{t-1},x_t]+b_i),
\]
\[
  o_t=\sigma(W_o[h_{t-1},x_t]+b_o),
\]
\[
  c_t=f_t\odot c_{t-1}+i_t\odot \tilde c_t.
\]

\subsection{Attention}

Scaled dot-product attention：
\[
  \operatorname{Attention}(Q,K,V)
  =
  \softmax\left(\frac{QK^T}{\sqrt{d_k}}\right)V.
\]
Transformer 把图像 patch 或 token 通过 self-attention 建模全局关系。ViT、DETR、CLIP、SAM 等现代视觉模型都来自这条线。

\section{Generative Models}

\subsection{Autoencoder}

Autoencoder 学习压缩表示：
\[
  z=f_\theta(x),\qquad \hat x=g_\phi(z),
\]
\[
  L=\|x-\hat x\|^2.
\]
瓶颈 $z$ 是对输入的低维表示。

\subsection{VAE}

VAE 学习概率潜变量：
\[
  q_\phi(z\mid x),\qquad p_\theta(x\mid z).
\]
优化 ELBO：
\[
  \mathcal L
  =
  \E_{q_\phi(z\mid x)}[\log p_\theta(x\mid z)]
  -
  \KL(q_\phi(z\mid x)\|p(z)).
\]

\subsection{GAN}

GAN 的 min-max 目标：
\[
  \min_G\max_D
  \E_{x\sim p_{\text{data}}}[\log D(x)]
  +
  \E_{z\sim p_z}[\log(1-D(G(z)))].
\]
生成模型让视觉从“识别世界”进一步走向“想象和生成世界”。

\section{Self-supervised Vision}

\subsection{Contrastive Learning}

对同一图像的不同增强视为正样本，对其他图像视为负样本。InfoNCE 形式：
\[
  \mathcal L_i=
  -\log
  \frac{\exp(\operatorname{sim}(z_i,z_i^+)/\tau)}
  {\exp(\operatorname{sim}(z_i,z_i^+)/\tau)+
  \sum_{j}\exp(\operatorname{sim}(z_i,z_j^-)/\tau)}.
\]
目标是让同一语义对象的表示靠近，不同对象远离。

\subsection{Masked Modeling}

Masked image modeling 随机遮挡 patch，让模型恢复缺失内容：
\[
  \hat x_{\mathcal M}=f_\theta(x_{\setminus \mathcal M}).
\]
这与空间智能里的世界模型很接近：系统要从部分观测推断完整世界。

\subsection{Vision-Language Representation}

CLIP 类模型把图像和文本投到同一空间。图像表示 $v$ 和文本表示 $t$ 的相似度：
\[
  s(v,t)=\frac{v^Tt}{\|v\|\|t\|}.
\]
这条线是空间智能从 perception 走向 language-grounded planning 的桥。

\section{CS231n 到空间智能的连接}

CS231n 给你神经视觉 backbone，CS231A 给你 3D 几何约束。现代空间智能通常是：
\[
  \text{image/video} \rightarrow \text{visual features}
  \rightarrow \text{geometry / depth / motion}
  \rightarrow \text{world state}
  \rightarrow \text{action}.
\]

学习顺序建议：
\begin{enumerate}
  \item 用 CS231n 学会视觉表示：CNN、Transformer、self-supervised features。
  \item 用 CS231A 学会几何：相机、对极、三角测量、SfM。
  \item 用 SLAM/三维重建把表示和几何合并：feature matching、pose estimation、BA、dense map。
  \item 用世界模型把视觉状态变成可预测、可决策的 latent state。
\end{enumerate}

\section*{公式速查}
\addcontentsline{toc}{section}{公式速查}

\begin{align*}
  &\text{kNN L2:} &
  d_2(x,x_i)&=\sqrt{\sum_p(x_p-x_{i,p})^2}.\\
  &\text{Linear classifier:} &
  s&=Wx+b.\\
  &\text{SVM:} &
  L_i&=\sum_{j\neq y_i}\max(0,s_j-s_{y_i}+\Delta).\\
  &\text{Softmax:} &
  p(y=j\mid x)&=\frac{e^{s_j}}{\sum_k e^{s_k}}.\\
  &\text{Cross entropy:} &
  L_i&=-\log p(y_i\mid x_i).\\
  &\text{SGD:} &
  \theta_{t+1}&=\theta_t-\eta\nabla_\theta L(\theta_t).\\
  &\text{ReLU:} &
  \ReLU(x)&=\max(0,x).\\
  &\text{Conv output:} &
  W_{\text{out}}&=\frac{W-F+2P}{S}+1.\\
  &\text{BatchNorm:} &
  \hat x_i&=\frac{x_i-\mu_B}{\sqrt{\sigma_B^2+\epsilon}}.\\
  &\text{ResNet:} &
  y&=F(x)+x.\\
  &\text{IoU:} &
  \operatorname{IoU}(A,B)&=\frac{|A\cap B|}{|A\cup B|}.\\
  &\text{Attention:} &
  \operatorname{Attention}(Q,K,V)&=\softmax\left(\frac{QK^T}{\sqrt{d_k}}\right)V.\\
  &\text{VAE ELBO:} &
  \mathcal L&=\E_{q_\phi}[\log p_\theta(x\mid z)]-\KL(q_\phi(z\mid x)\|p(z)).\\
  &\text{GAN:} &
  \min_G\max_D&\ \E_x[\log D(x)]+\E_z[\log(1-D(G(z)))].
\end{align*}

\section*{GoodNotes 页码建议}
\addcontentsline{toc}{section}{GoodNotes 页码建议}

\begin{center}
\begin{tabular}{p{0.12\linewidth}p{0.34\linewidth}p{0.44\linewidth}}
\toprule
页码 & 标题 & 必画图 \\
\midrule
001 & Image Classification & train/val/test 和 kNN 投票。 \\
002 & SVM and Softmax & 分数向量、margin、概率。 \\
003 & Backprop & 计算图和局部梯度。 \\
004 & Neural Network & affine-ReLU-affine 结构。 \\
005 & CNN & 卷积核滑动、输出尺寸、参数量。 \\
006 & Training Tricks & BatchNorm、Dropout、Momentum。 \\
007 & Detection/Segmentation & box、IoU、mask。 \\
008 & Attention & $Q,K,V$ 和 token 关系图。 \\
009 & Generative Models & encoder/decoder、VAE、GAN。 \\
010 & Spatial Bridge & CS231n feature 到 CS231A geometry 的连接图。 \\
\bottomrule
\end{tabular}
\end{center}

\end{document}
